\section{Introduction}
\label{intro}

Multi-agent systems (MAS) are increasingly proposed as a substrate for problems that exceed the practical capacity of single-agent learners, such as disaster response, transportation, and large-scale web automation, where parallelism, partial observability, and large scale dynamics are intrinsic to the task \cite{yuan2023survey, wong2023deep}. Multi-Agent Reinforcement Learning (MARL) offers a framework for such systems to learn cooperative joint policies through environmental interaction \cite{yuan2023survey, oroojlooy2023review}, and recent work extends this paradigm to LLM-based agents that integrate planning, memory, and natural-language communication into goal-directed behaviour \cite{xi2025rise, sun2024llm}.

Yet across both lines of work, agents typically treat one another as transient sources of information rather than as enduring partners. Coordination is achieved through learned vector messages \cite{zhu2024survey}, LLM-mediated dialogue \cite{sun2024llm}, or attention over the population \cite{iqbal2019actor}, all recomputed at every step and dissolving when the interaction ends. Self-evolving frameworks soften this by letting agents adapt internal components online \cite{fang2025comprehensive, gao2025survey}, but the population structure itself remains fixed or implicit. The limitation is most consequential in \emph{open environments} where population and tasks shift dynamically \cite{yuan2023survey}, settings that demand persistent social relationships rather than just better per-step coordination \cite{licua2025mindforge}.

Human intelligence offers a model for what is missing. Humans learn not only from direct experience but from one another, and this social learning is what allows knowledge to accumulate across a group rather than being rediscovered by each individual \cite{tomasello1993cultural, hoppitt2013social, boyd2011cultural}. Moreover, social learning is selective: it relies on persistent relationships that determine whom to learn from and what to retain \cite{kendal2018social}. Existing work on social learning between artificial agents \cite{ndousse2021emergent, hillier2024social} captures the act of learning from others but not the persistent bonds that determine whom to learn from over time.

We argue that the missing component in MAS is \emph{structural plasticity at the population level}. Just as the brain's capacity to learn arises not from the intelligence of individual neurons but from the dynamic plasticity of the synapses between them \cite{hebb2005organization, kusmierz2017learning}, a MAS can be viewed as an adaptive network in which inter-agent bonds form the substrate of social intelligence. Under this view, agents play the role of neurons, social bonds play the role of synapses, and reward feedback plays the role of the modulatory signal. We propose that applying a reward-modulated Hebbian rule to inter-agent bonds produces stable social assemblies, meaning pairs and groups of agents that preferentially collaborate, and provides a structural basis for credit propagation, and targeted communication in MAS.

\paragraph{Contributions.} This research makes the following contributions: (i) we formalise a \emph{neuron to agent isomorphism} that maps a MAS to a time-evolving weighted graph governed by a reward-modulated Hebbian rule; (ii) we operationalise the graph through \emph{reward diffusion} across bonded agents and adaptation of partner's beliefs; (iii) we release \emph{Five Chambers}, an open-source Craftium environment with five sequential rooms that require progressively tighter coordination, from solo learning to joint combat; and (iv) we evaluate the framework across a layered comparison of LLM-only agents, LLM agents augmented with a MAPPO RL layer, and the addition of Hebbian Social bonds to the system.
\todo{Here I need to decide what experiments will be added in the end and if I should change this}

\paragraph{Research questions.} We organise our work around four research questions:
\begin{quote}
\textbf{RQ1} Does reward-modulated Hebbian plasticity over inter-agent bonds improve cooperative task performance, compared to LLM-only and RL baselines without a social graph?

\textbf{RQ2} How does the social graph topology evolve over training, and do emergent bond patterns reflect the cooperative structure of the environment?

\textbf{RQ3} Does reward diffusion across the Hebbian graph propagate learning signals across the team, accelerating skill acquisition compared to agents that receive only their own reward?
\end{quote}
